\clearpage
\SourceLeaf{451}
\begin{center}\large [ABWICKLUNGSFÄHIGE FLÄCHEN.]\end{center}
\begin{center}[1.]\end{center}
\noindent\hspace*{1.2em}1825, 4. Dezember.\quad Das Nachdenken über die Theorie der abwicklungsfähigen Flächen hat uns folgendes elegante analytische Theorem an die Hand gegeben.

Ist $z$ eine solche Funktion von $x,y$, dass
\[
\frac{\partial\partial z}{\partial x^2}.\frac{\partial\partial z}{\partial y^2}
=\left(\frac{\partial\partial z}{\partial x\partial y}\right)^2,
\]
\noindent so lassen sich statt $x,y$ zwei andere veränderliche Größen $t,u$ einführen, sodass $x,y,z$, als Funktionen von $t,u$ ausgedrückt, in Bezug auf $u$ lineare Funktionen sind. In der Tat kann man für $u$ jede beliebige lineare Funktion von $x,y,z$ wählen.

\Needspace{7\baselineskip}
\begin{center}[2.]\end{center}
\begin{center}\textsc{Gauss} an \textsc{Olbers}. Göttingen, Juli 1828.\end{center}
\noindent\hspace*{1.2em}. . . . . Die Übersetzung des Artikels der Göttingischen Gelehrten Anzeigen über meine Abhandlung zu den krummen Flächen stammt, wie mir \textsc{Schumacher} sagte, von \textsc{Baily}. Er erzählte mir zugleich von dem Angriff des p.\ \textsc{Fayolles} und wünschte etwas, das \textsc{Baily} zur Erwiderung verwenden könne,
